\subsection{Review-Frame Discipline and Scope Contract}\label{app:review_frame}

This appendix makes the author's reviewing stance explicit. It accompanies the mechanical certificate (\appref{app:cert_scope}) and is intended to reduce the perception that pre-work pre-empts review. The appendix is exhaustive about framing discipline. It is not a substitute for the body claims, the proofs, or the empirical measurement.

\paragraph{Author's stance.}
We did mechanical pre-work. The certificate suite covers a finite class of objections by precompiling deterministic checks against them. Mechanical work narrows one surface. It does not narrow the field of valid reviewer objection. Where a check has been precompiled, a check-level failure can be pointed at directly. Where no check has been precompiled, the certificate is silent and the reviewer's judgment stands on its own terms. The author's reading of an empirical pattern is offered, not imposed. A reviewer who reads the same pattern differently is correct to record that reading.

\paragraph{Reviewer-frame acknowledgment.}
Different reviewers read the paper from different operating frames. A theorist reviewer may read for proof structure, assumption discipline, and tightness of the bound. An experimentalist reviewer may read for protocol, sample size, and reproducibility. A methodologist may read for the shape of the audit observer and whether the deterministic substrate is what it claims to be. Each frame is legitimate. The paper carries pointers for each. We do not claim to know which frame the assigned reviewer occupies, and we do not encode that claim into the prose.

\paragraph{What the certificate covers, in plain language.}
The certificate suite (\appref{app:cert_scope}) organizes into six suites. The claim-text suite verifies that the paper's numeric and structural claims appear in the paper without drift. The data-ties suite verifies that those claims are tied to the source data. The artifact-lineage suite verifies that figures bind to scripts and data. The provenance suite verifies that the runtime environment is recorded and reproducible. The statistical-hygiene suite verifies that caveats are tracked and uncertainty surfaces are not concealed. The submission-hygiene suite verifies that citations, links, license, and PDF readiness are clean. Three additional layers cover the audit observer substrate, with sources at \texttt{ci/audit/} and the canonicality ledger at \texttt{ci/audit/CANONICALITY\_LEDGER.md}. None of these suites adjudicate scientific truth. They adjudicate the form of the artifact.

\paragraph{Substantive disagreement versus check-level disagreement.}
A check-level disagreement is one that maps onto a precompiled check. A reviewer who finds such a disagreement can name the check, point to its output, and ask the authors to address the failure in correspondence. A substantive disagreement is one outside the precompiled coverage. It may concern the framing of the bound, the choice of empirical methodology, the interpretation of the audit observer's verdict, the suitability of the displacement contract, or any aspect of the work that the certificate does not adjudicate. We do not treat substantive disagreement as inferior to check-level disagreement. The certificate covers the surface. The reading covers the meaning.

\paragraph{Scope contract.}
The certified domain is explicit, verifiable multi-constraint tasks with deterministic validators. Within that domain, the load-bearing claims hold under the framework. The regime index orders failure regimes ($r_s = 1.0$). The smooth-pivot decomposition under audit-observer measurement gives the reported pooled fractions. The high-tier refutation count is zero across the 5{,}472-trial measurement run. The pre-generation routing achieves $1.8\%$ regret on calibration. These are the empirical surfaces the certificate adjudicates.

The non-certified domain includes implicit residual judgment, broad human preference alignment, semantic qualia, and unconstrained natural-language helpfulness. Where the paper discusses these, it does so as scope description rather than as certified result. A reviewer who reads the paper as overclaiming is asked to identify a specific load-bearing claim that crosses the certified-domain boundary. The certificate lists every load-bearing claim by identifier in \texttt{ci/claim\_data\_ties.json}.

The bridge between certified and non-certified is decomposition. Implicit constraints may be decomposed into explicit verifiable subconstraints plus a residual judgment layer. Only the explicit layer is certified. The decomposition itself is described in Algorithm~\ref{alg:routing} and is not claimed to be complete.

\paragraph{Judge-free verdict layer with explicit callouts.}
LLM-as-judge is itself the research subject of this paper. We study how LLMs behave under multi-constraint judging. Within that frame, the cert system's verdict layer is deterministic so that the cert does not recursively rely on the behavior under study. Core empirical pass-fail outcomes are decided by AST-checked Python and regex format verifiers. The audit observer renders its verdicts via a deterministic predicate ladder over verifier outcomes. No LLM is in the decision loop for any cert layer's pass-fail determination. The L30 audit-observer purity check verifies this by enumerating every import in every substrate file and confirming none are LLM SDKs. The substrate's import graph is mechanically checkable. We acknowledge LLM appearances elsewhere in the broader pipeline. The implicit-k constraint decompression pipeline uses an LLM as a tagger that selects from a fixed 34-rule deterministic verifier library. The image-format Pass B pilot is hand-rated by a single rater. Future paper-surface fault-injection work may use LLM-as-judge as part of its oracle. None of these are concealed cert dependencies. They are acknowledged adjacent methodology, distinct from the deterministic verdict layer.

\paragraph{Structured residue and successor questions.}
The framework predicts that collapse leaves structured residue. The residue includes staging traces, refusals, the residual smooth-regime trials, and the verifier-surface gap between model families flagged by the audit observer. We identify the question of how this residue integrates into a new feasible regime as a structural successor to the present work. We do not claim it is the only successor question. A reviewer in a different frame may identify a different successor as the natural next theoretical object. The structured-residue framing is offered to make our own reading explicit, not to constrain other readings. We do not solve the residue-integration question here.

\paragraph{Verifier-surface mismatch.}
The audit observer's stream-level decision flagged a verifier-surface mismatch between the Claude family and the open-weight family. High-tier pass rates differ by approximately a factor of two ($83\%$ for Claude family, $45\%$ for open-weight family). The framework's H\_B3 verdict (zero true-certificate refutations) survives across both families and remains the load-bearing empirical claim. The H\_B1 and H\_B2 verdicts pass under per-(observer, model) calibration. We disclose the mismatch rather than concealing it. The mismatch may be a genuine family-level difference in alignment under the displacement contract. It may be an artifact of the specific verifier rules used in this study. It may be both. We do not adjudicate. The substrate, the run manifest, and the per-packet provenance are released so that a reviewer or replicator can examine the gap directly.

\paragraph{Forward direction: paper-surface fault-injection.}
The structured-residue framing above points toward a methodological direction we anticipate developing in follow-on work. A scientific manuscript can be treated as an observer-facing artifact whose claims, figures, and certificate interfaces should remain stable under bounded perturbation. We intend to formalize a fault-injection protocol for manuscripts (qualifier deletion, hostile paraphrase, theorem-assumption erasure, frame-swap, fragment-only reading) and report the resulting fragility map as a separate diagnostic artifact, distinct from the present submission's evidence base. As an early signal we ran a small static analysis of representative perturbations against this paper's existing cert chain. Numeric and table-cell mutations were caught at $100\%$. Lexicon swaps, qualifier deletions, citation drops, and hostile paraphrases survived at $100\%$. The aggregate survival rate was $47.1\%$ across $17$ applicable mutations. The deterministic cert chain protects the numeric and data-tied surfaces well. The prose-semantic surface remains unmonitored. This is the empirical opening for the follow-on program. The protocol, harness, and full fragility map are not part of this submission.

\paragraph{What this stance achieves and what it does not.}
This appendix does not eliminate disagreement. It does not substitute for review. It does not pre-empt the AC's meta-review. It only narrows the surface on which the authors have done deterministic pre-work and discloses, in long form, the boundary between that surface and the field of valid reviewer judgment that lies outside it. The authors prefer this disclosure to the alternative of an unstated scope contract that a reviewer must reverse-engineer from the prose.
